\documentclass{article}
\usepackage{amsmath, amssymb, amsthm}
\usepackage{hyperref}
\usepackage{geometry}
\geometry{margin=1in}

\title{Erd\H{o}s Problem \#451}
\date{Last updated: 2025-08-31}
\author{Source: erdosproblems.com}

\begin{document}
\maketitle

\noindent\textbf{Status:} OPEN \quad \textbf{No prize}

\noindent\textbf{Tags:} number theory

\medskip
\noindent\textbf{Problem Statement:}

Estimate $n_k$, the smallest integer $>2k$ such that $\prod_{1\leq i\leq k}(n_k-i)$ has no prime factor in $(k,2k)$.

\bigskip
\noindent\textbf{Remarks:}

Erd\H{o}s and Graham write 'we can prove $n_k>k^{1+c}$ but no doubt much more is true'. 

In \cite{Er79d} Erd\H{o}s writes that probably $n_k<e^{o(k)}$ but $n_k>k^d$ for all constant $d$.

Adenwalla observes that an easy upper bound is $n_k\leq \prod_{k<p<2k}p=e^{O(k)}$.

van Doorn and Tang \cite{vDTa26} (using an argument generated by GPT 5.5 Pro and Tang) have proved that\[n_k> \exp\left(c\frac{(\log k)^2}{\log\log k}\right)\]for some constant $c>0$. 

See also [1095].

\bigskip
\noindent\textbf{References:}

[Er79d] Erd\H{o}s, P., Some unconventional problems in number theory. Acta Math. Acad. Sci. Hungar. (1979), 71-80.
 
 [vDTa26] W. van Doorn and Q. Tang, Consecutive integers free of certain prime factors. arXiv:2606.19863 (2026).

\bigskip
\noindent\small{Source: \url{https://www.erdosproblems.com/451}}
\end{document}
